\documentclass[11pt,a4paper]{article}

\usepackage[utf8]{inputenc}
\usepackage[T1]{fontenc}
\usepackage[spanish,es-noquoting,es-nodecimaldot,es-tabla]{babel}
\usepackage{amsmath,amssymb}
\usepackage{siunitx}
\sisetup{detect-all, per-mode=symbol}
\usepackage{booktabs}
\usepackage{array}
\usepackage[section]{placeins}
\newcolumntype{L}[1]{>{\raggedright\arraybackslash}p{#1}}
\usepackage{microtype}
\setlength{\emergencystretch}{2em}
\usepackage[margin=2.6cm]{geometry}
\usepackage{xcolor}
\usepackage[colorlinks=true,linkcolor=azulosc,citecolor=azulosc,urlcolor=azulosc]{hyperref}

\definecolor{azulosc}{HTML}{1F4E79}
\definecolor{grisclaro}{HTML}{F2F2F2}

\usepackage{tcolorbox}
\tcbuselibrary{skins,breakable}
\newtcolorbox{nota}[1][]{colback=grisclaro, colframe=azulosc, boxrule=0.6pt,
  arc=2pt, left=6pt, right=6pt, top=5pt, bottom=5pt, breakable, #1}

\title{Batería de pruebas para el régimen de autogravedad\\
       \large Diseño, justificación de cada decisión y protocolo de medida}
\author{Notas de trabajo}
\date{23 de septiembre de 2026}

\begin{document}
\maketitle

\begin{abstract}
Se propone una batería de 28 simulaciones para decidir si, en la reducción de
momento angular fijo $f = F(r,p_r)\,\delta(L-L_0)$ con fondo isócrono, la
autogravedad de la componente autogravitante puede sostener una oscilación que
no se amortigua, o si el resultado es siempre amortiguamiento por mezcla de
fases. El argumento central del diseño es que \textbf{la masa del dato inicial
no es el parámetro de control}: lo es el cociente $\eta$ entre el corrimiento
de frecuencias que induce la autogravedad y el ancho de la banda de frecuencias
orbitales de la propia distribución. Las corridas usan una Maxwelliana rebajada
(el modelo de Wilson): es el análogo en energía de Maxwell--Boltzmann, truncado
por marea, de soporte compacto y banda de frecuencias exacta, y monótona, de
modo que por el teorema de Antonov no admite modos radiales crecientes. Todas
las cifras del documento están calculadas con los guiones del repositorio
(\texttt{reproducir/scripts/eta.py} y \texttt{colas\_\allowbreak libres.py}), no
estimadas. El costo total es de aproximadamente una hora y media de cómputo con
cuatro hilos.
\end{abstract}

\section{Resumen ejecutivo}

\begin{nota}
\textbf{Qué se prueba.} Que el comportamiento tardío de una perturbación
(amortiguarse o seguir oscilando) está controlado por
\[
  \eta \;=\; \frac{\text{corrimiento de la frecuencia media por autogravedad}}
                   {\text{ancho de la banda de frecuencias orbitales}},
\]
y no por la masa $a_0$ por separado ni por el ancho $\sigma_J$ por separado.

\textbf{Cómo se prueba.} Barriendo $\eta$ de $0.1$ a $1.5$ (bloque~A), y
llegando al \emph{mismo} $\eta$ por caminos distintos: tres anchos del soporte
$J_t$ y dos valores de $L_0$ (bloques~A y~D). Si el resultado colapsa en
$\eta$, el mecanismo queda demostrado.

\textbf{Qué más se obtiene.} La frontera del régimen lineal (bloque~C), la
convergencia numérica en las cuatro fuentes de error independientes
(bloque~B) y tres controles de familia (bloque~E): una distribución hueca en
$J=0$, la gaussiana ya validada y el límite térmico de la familia principal.
\end{nota}

\begin{table}[htbp]
\centering
\begin{tabular}{llrr}
\toprule
Bloque & Propósito & Corridas & Costo \\
\midrule
A & Barrido en $\eta$ y prueba de degeneración & 9 & \SI{23}{min}\\
B & Convergencia: $N$, $\Delta r$, $\Delta t$, $N_{pc}$ & 7 & \SI{37}{min}\\
C & Linealidad y atrapamiento (eje $\mu$) & 7 & \SI{27}{min}\\
D & Control con otro momento angular $L_0$ & 2 & \SI{3}{min}\\
E & Controles de familia & 3 & \SI{9}{min}\\
\midrule
 & \textbf{Total} & \textbf{28} & $\approx$\,\SI{1.6}{h}\\
\bottomrule
\end{tabular}
\caption{Estructura de la batería. Los tiempos son con cuatro hilos, calibrados
con una corrida real: \SI{100}{s} para $10^4$ partículas, $6\times10^4$ pasos,
yoshida4 y autogravedad.}
\end{table}

\section{Marco físico}

Esta sección construye, paso a paso, las tres cantidades que organizan la
batería: la banda de frecuencias orbitales, el corrimiento que produce en ella
la autogravedad y la amplitud de la perturbación del potencial. Cada paso usa
una aproximación; para cada una se dice cuál es, por qué es razonable y cómo se
comprobó con cálculos o corridas reales. La tabla~\ref{tab:aprox} al final de la
sección las resume.

\subsection{La reducción a momento angular fijo}

En simetría esférica, la distribución depende de $(r,p_r,L)$. La reducción que
resuelve el código es
\[
  f(r,p_r,L) = F(r,p_r)\,\delta(L-L_0),
\]
es decir, todas las partículas tienen el mismo momento angular. El movimiento
radial de cada una es entonces unidimensional, con hamiltoniano
\[
  H = \tfrac12 p_r^2 + \Phi_{\rm ef}(r), \qquad
  \Phi_{\rm ef}(r) = \Phi(r) + \frac{L_0^2}{2r^2},
\]
donde $\Phi = \Phi_{\rm iso} + \Phi_{\rm self}$ es el potencial total: el del
fondo isócrono, $\Phi_{\rm iso} = -1/(1+\sqrt{1+r^2})$ con masa y escala unidad,
más el de la propia distribución. La medida del espacio fase que deja la delta
es $8\pi^2 L_0\,dr\,dp_r$, de modo que la masa es
$M = 8\pi^2L_0\int F\,dr\,dp_r$.

Esta reducción no es una aproximación \emph{dentro} del modelo, pero sí es un
modelo extremadamente anisótropo. Todo lo que sigue vale para él, y la
comparación con resultados de sistemas isótropos (como el de Hadžić y
colaboradores) debe hacerse con ese cuidado.

\subsection{Variables acción-ángulo y frecuencias}

Como $H$ no depende del tiempo y el movimiento es acotado, existen variables
acción-ángulo. La acción radial es el área de la órbita en el plano
$(r,p_r)$ dividida por $2\pi$,
\[
  J = \frac{1}{\pi}\int_{r_-}^{r_+} \sqrt{2\,(E-\Phi_{\rm ef}(r))}\;dr ,
\]
con $r_\pm$ los radios de retorno, y el ángulo $Q$ avanza uniformemente,
$\dot Q = \Omega(J)$, con
\[
  \Omega(J) = \frac{dE}{dJ} = \frac{2\pi}{T(J)}, \qquad
  T = 2\int_{r_-}^{r_+}\frac{dr}{\sqrt{2(E-\Phi_{\rm ef})}} .
\]
Solo el isócrono tiene forma cerrada:
\[
  E(J) = -\frac{1}{2\,(J+c)^2}, \qquad
  \Omega(J) = \frac{1}{(J+c)^3}, \qquad
  \frac{d\Omega}{dJ} = -\frac{3}{(J+c)^4}, \qquad
  c = \tfrac12\left(L_0+\sqrt{L_0^2+4}\right).
\]
Dos consecuencias se usan en todo el documento. La frecuencia \emph{decrece}
con la acción: las órbitas casi circulares ($J\to0$) son las más rápidas, y
$\Omega_{\max} = \Omega(0) = c^{-3}$. Y la derivada no se anula, así que el mapa
$J\mapsto\Omega$ es invertible sobre todo el soporte.

Con autogravedad no hay forma cerrada. El mapa se calcula numéricamente
(\texttt{aa\_\allowbreak numerico.py}): para cada energía se buscan los radios de retorno por
bisección y las integrales se hacen con el cambio $r = r_m + r_a\sin\theta$, que
elimina la singularidad de raíz cuadrada en los puntos de retorno, y cuadratura
de Gauss--Legendre. $E(J)$ queda tabulada y $\Omega = dE/dJ$ se obtiene por
diferencias finitas. Así calcula \texttt{eta.py} todas las bandas de este
documento.

\subsection{La banda de frecuencias}

El soporte de $F_{\rm eq}(J)$ ocupa un intervalo de acciones y, a través de
$\Omega(J)$, una banda de frecuencias $[\Omega_{\min},\Omega_{\max}]$. Esa banda es
el \emph{espectro esencial} del problema linealizado: el conjunto de
frecuencias con las que las partículas pueden resonar. Su ancho y su posición
son los protagonistas del resto del documento.

\paragraph{Aproximación 1: dónde termina el soporte.} Las distribuciones de la
batería (sección~\ref{sec:familia}) tienen soporte compacto $[0,J_t]$: la banda
es exactamente $[\Omega(J_t),\,\Omega(0)]$ y no hay nada que elegir. La
gaussiana $J^2e^{-J^2/\sigma_J^2}$ de las corridas de Landau y de la corrida puente
E3 no tiene borde, así que hay que fijar un umbral: su soporte es donde
$F_{\rm eq}\ge 10^{-2}\,F_{\max}$. Es una convención: con otro umbral cambian los
bordes y, con ellos, el valor absoluto de $\eta$. El borde de la familia
compacta se eligió en el extremo superior de ese soporte, $J_t = 2.763\,\sigma_J$,
para que las dos familias tengan la misma banda.

\paragraph{Aproximación 2: el ancho relativo a primer orden.} Si el intervalo de
acciones $\Delta J$ es chico frente a $J+c$, un desarrollo de Taylor da
\[
  \frac{\Delta\Omega}{\Omega} \;\approx\; \frac{|d\Omega/dJ|\,\Delta J}{\Omega}
  \;=\; \frac{3\,\Delta J}{J+c}.
\]
Esta fórmula explica de un vistazo las dos perillas: el ancho crece con
$\sigma_J$ (que fija $\Delta J$) y decrece con $L_0$ (que fija $c$). Se comprobó
contra el cálculo exacto del isócrono:

\begin{table}[htbp]
\centering
\small
\begin{tabular}{llcccc}
\toprule
soporte & ancho & soporte en $J$ & banda $[\Omega_{\min},\Omega_{\max}]$ &
  $\Delta\Omega/\Omega$ exacto & $3\Delta J/(J+c)$\\
\midrule
gaussiana & $\sigma_J=0.10$ & $[0.0061,\,0.2763]$ & $[0.0513,\,0.0705]$ & 0.315 & 0.317\\
          & $\sigma_J=0.05$ & $[0.0031,\,0.1382]$ & $[0.0601,\,0.0708]$ & 0.163 & 0.163\\
          & $\sigma_J=0.03$ & $[0.0019,\,0.0829]$ & $[0.0642,\,0.0709]$ & 0.099 & 0.099\\
\midrule
compacto  & $J_t=0.276$ & $[0,\,0.276]$ & $[0.0514,\,0.0711]$ & 0.322 & 0.324\\
          & $J_t=0.138$ & $[0,\,0.138]$ & $[0.0602,\,0.0711]$ & 0.166 & 0.167\\
          & $J_t=0.083$ & $[0,\,0.083]$ & $[0.0642,\,0.0711]$ & 0.101 & 0.101\\
\bottomrule
\end{tabular}
\caption{Bandas del isócrono con $L_0=2$ ($c=2.414$), sin autogravedad;
$\Omega$ es el centro de la banda. La aproximación lineal acierta a menos del
\SI{1}{\%}; el error crece, como debe, con el ancho del intervalo. Sin masa, la
banda de soporte compacto es la misma para todas las familias; llega hasta
$\Omega(0)=c^{-3}$ porque no hay umbral que recorte las órbitas casi
circulares.}
\label{tab:bandas}
\end{table}

\subsection{Mezcla de fases: por qué decae $h_k$ aun sin autogravedad}

Sin autogravedad (o despreciándola), la evolución en variables acción-ángulo es
trivial: cada punto conserva su acción y su ángulo gira,
\[
  F(Q,J,t) = F\big(Q-\Omega(J)\,t,\;J,\;0\big).
\]
Con el dato inicial $F = F_{\rm eq}(J)\,[1+\varepsilon\cos Q]$, el coeficiente de
Fourier $k=1$ es $\tfrac{\varepsilon}{2}F_{\rm eq}(J)\,e^{-i\Omega(J)t}$, y el
observable, pesado por la función de prueba $B(J)$, vale
\[
  h_1(t) \;\propto\; \int dJ\;B(J)\,F_{\rm eq}(J)\,e^{-i\Omega(J)t}
  \;=\; \int d\Omega\; w(\Omega)\,e^{-i\Omega t},
  \qquad
  w(\Omega) = \frac{B\,F_{\rm eq}}{|d\Omega/dJ|}.
\]
El segundo paso es solo el cambio de variable $J\to\Omega$, legítimo porque
$\Omega(J)$ es monótona. Resultado: \textbf{$h_1(t)$ es la transformada de Fourier
de la distribución de frecuencias $w(\Omega)$}. Por el lema de
Riemann--Lebesgue tiende a cero, y la \emph{forma} en que lo hace la dicta la
regularidad de $w$:

\begin{itemize}
\item Si $w$ es suave y localizada, el decaimiento es más rápido que cualquier
potencia. Para una $w$ aproximadamente gaussiana de desviación $\sigma_\Omega$,
$|h_1(t)|\approx|h_1(0)|\,e^{-\sigma_\Omega^2t^2/2}$: decaimiento gaussiano en el
tiempo.
\item Si $w$ tiene un borde donde se anula como $(\Omega-\Omega_{\min})^k$, ese
borde aporta una cola algebraica $\propto t^{-(k+1)}$ (la asintótica estándar de
un extremo de integración). Cuanto más brusco el borde, más lenta la cola. Esta
es la conexión intuitiva con la sección~\ref{sec:borde}.
\end{itemize}

La escala de tiempo es la inversa del ancho. Definimos
\[
  \tau_k = \frac{2\pi}{k\,\Delta\Omega},
\]
con $\Delta\Omega$ el ancho del soporte de la tabla~\ref{tab:bandas}.

\paragraph{Aproximación 3: $\tau_k$ es una escala, no una ley.} El decaimiento
real no es exponencial con tasa $1/\tau_k$. Lo que $\tau_k$ da es el orden de
magnitud del tiempo en que la mezcla es completa. La relación con la
desviación $\sigma_\Omega$ depende del umbral del soporte: para un perfil
gaussiano, el punto donde cae al \SI{1}{\%} está a $3.03\,\sigma$ del centro, así
que $\Delta\Omega\approx 6.1\,\sigma_\Omega$.

\paragraph{Comprobación con las corridas de Landau.} La corrida
\texttt{landau\_\allowbreak libre} evoluciona la perturbación sin autogravedad en el
potencial del equilibrio con $\sigma_J=0.10$. Si el decaimiento es gaussiano,
$\ln(|h_1(0)|/|h_1(t)|) = \sigma_\Omega^2t^2/2$ con $\sigma_\Omega$ constante.
Los datos dan

\begin{table}[htbp]
\centering
\begin{tabular}{rccc}
\toprule
$t$ & $|h_1|/|h_1(0)|$ & $\sigma_\Omega$ inferido & $\Delta\Omega/\sigma_\Omega$\\
\midrule
400 & $5.5\times10^{-1}$ & $2.74\times10^{-3}$ & 6.97\\
800 & $8.0\times10^{-2}$ & $2.81\times10^{-3}$ & 6.80\\
1000 & $1.6\times10^{-2}$ & $2.88\times10^{-3}$ & 6.63\\
1200 & $1.5\times10^{-3}$ & $3.00\times10^{-3}$ & 6.36\\
1400 & $9.9\times10^{-5}$ & $3.07\times10^{-3}$ & 6.23\\
\bottomrule
\end{tabular}
\caption{Decaimiento libre medido. $\sigma_\Omega$ varía solo un \SI{12}{\%} mientras
$|h_1|$ cae cuatro órdenes de magnitud: el decaimiento es casi exactamente
gaussiano, como predice la transformada de Fourier de una $w$ suave. El ligero
aumento tardío es la cola no gaussiana de $w$. El cociente
$\Delta\Omega/\sigma_\Omega\approx 6.2$--$7.0$ coincide con el $6.1$ esperado.}
\end{table}

Una consecuencia práctica: la parte libre de $h_1$ alcanza el piso numérico
hacia $t\approx 5\,\tau_1$. De ahí la duración $t_{\rm fin}=10\,\tau_1$ de la
batería: la mitad de la corrida para que termine la mezcla libre y la otra mitad
para ajustar lo que quede. Para $\delta\Phi$ la cuenta es otra, y es el tema del
párrafo siguiente.

\paragraph{La cola libre de $\delta\Phi$ y la forma de $F_{\rm eq}$ en $J=0$.}
$h_1$ pesa con $B(J)$, que es suave y se anula como $J^2$. El observable principal
del protocolo (sección~\ref{sec:protocolo}) es en cambio la perturbación del
potencial en un radio fijo, que sin autogravedad también es una cuadratura
exacta:
\[
  \delta\Phi(r,t) = 8\pi^2L_0\,\varepsilon\;
  \mathrm{Re}\!\int dJ\;F_{\rm eq}(J)\,s(J)\,C(r,J)\,e^{-i\Omega(J)t},
  \qquad
  C(r,J) = \int dQ\;e^{iQ}\,G\big(r,\,r(Q,J)\big),
\]
con $G(r,r')=-1/\max(r,r')$ la función de Green de capas esféricas y $s(J)$ el
factor radial de la perturbación ($s=1$ en 11\_landau; sección~\ref{sec:perfil}).
Dos propiedades de $C$ fijan la cola, por la regla de la sección anterior (un
integrando que va como $x^a$ en un punto aporta $t^{-(a+1)}$):
\begin{itemize}
\item Cerca de la órbita circular, $r(Q,J)\approx r_c+\sqrt{2J/\kappa}\,\cos Q$, así
que $C\propto\sqrt J$. El extremo $J=0$ de la banda ($\Omega_{\max}$) aporta un
integrando $\propto F_{\rm eq}\,s\,J^{1/2}$. Si $F_{\rm eq}(0)>0$, como en toda
distribución monótona, la cola es $t^{-3/2}$ con $s=1$, $t^{-2}$ con
$s=\sqrt{J/J_t}$ y $t^{-3}$ con $s=(J/J_t)^{3/2}$. Si $F_{\rm eq}\propto J^2$
el integrando va como $J^{5/2}$, y la cola es $t^{-7/2}$.
\item Donde un punto de retorno de la órbita cruza el radio $r$, la fracción de
órbita más allá de $r$ crece como $(J-J_*)^{3/2}$, y $C$ hereda esa
singularidad: una cola $t^{-5/2}$ con frecuencia $\Omega(J_*)$, dentro de la
banda, \emph{para cualquier} $F_{\rm eq}$. $h_1$ no la tiene porque no mira un
radio fijo.
\end{itemize}
Así que la parte libre de $\delta\Phi$ decae algebraicamente, aun con la
gaussiana. La tabla~\ref{tab:colas} da los valores exactos para las familias de
la sección~\ref{sec:familia}.

\begin{table}[htbp]
\centering
\small
\setlength{\tabcolsep}{3.6pt}
\begin{tabular}{llcccccc}
\toprule
$F_{\rm eq}$ & $s(J)$ & $|h_1|/|h_1(0)|$ &
  \multicolumn{3}{c}{$\|\delta\Phi\|/\|\delta\Phi(0)\|$} & pendiente &
  $\|\delta\Phi(0)\|/M$\\
 & & $10\tau_1$ & $2\tau_1$ & $5\tau_1$ & $10\tau_1$ & $20$--$40\,\tau_1$ & \\
\midrule
gaussiana & 1 &
  $1.7\times10^{-6}$ & $7.9\times10^{-2}$ & $1.2\times10^{-3}$ & $1.2\times10^{-4}$ & $-2.4$ & 1.00\\
$J^2(J_t-J)^2$ & 1 &
  $9.7\times10^{-5}$ & $4.3\times10^{-2}$ & $3.1\times10^{-3}$ & $4.2\times10^{-4}$ & $-3.0$ & 1.09\\
Wilson, $W_0=3$ & 1 &
  $1.7\times10^{-4}$ & $1.9\times10^{-1}$ & $5.0\times10^{-2}$ & $1.8\times10^{-2}$ & $-1.5$ & 0.66\\
Wilson, $W_0=3$ & $\sqrt{J/J_t}$ &
  $5.6\times10^{-5}$ & $1.2\times10^{-1}$ & $2.0\times10^{-2}$ & $4.9\times10^{-3}$ & $-2.0$ & 0.32\\
Wilson, $W_0=3$ & $(J/J_t)^{3/2}$ &
  $4.3\times10^{-5}$ & $4.4\times10^{-2}$ & $2.4\times10^{-3}$ & $3.0\times10^{-4}$ & $-2.8$ & 0.11\\
Wilson, $W_0=6$ & $(J/J_t)^{3/2}$ &
  $3.6\times10^{-5}$ & $1.0\times10^{-1}$ & $7.7\times10^{-3}$ & $1.0\times10^{-3}$ & $-2.9$ & 0.06\\
\bottomrule
\end{tabular}
\caption{Decaimiento libre exacto (\texttt{colas\_\allowbreak libres.py}) con el
ancho de A4, $\sigma_J=0.05$ o $J_t=0.138$: isócrono desnudo con $L_0=2$, $B(J)$
con $j_1=s_{j1}=0.05$ y $\|\delta\Phi\|$ el valor rms en $r\in[3,15]$. ``Wilson''
es la Maxwelliana rebajada con $g=2$ de la sección~\ref{sec:familia}. La
pendiente es $d\ln\|\delta\Phi\|/d\ln t$ entre $20$ y $40\,\tau_1$ y confirma los
exponentes predichos: $-3/2$, $-2$ y, con $s=(J/J_t)^{3/2}$, una cola del borde
$J=0$ más rápida que la de los puntos de retorno ($-5/2$). La última columna es
$\delta\Phi(0)$ por unidad de masa y de $\varepsilon$, relativa a la gaussiana:
cuánto potencial produce la perturbación. La cuadratura reproduce a dos cifras
el decaimiento de \texttt{landau\_\allowbreak libre} (la suma sobre las
partículas de 11\_landau) en $t=400$--$1400$, y cambia menos de
$3\times10^{-4}$ (relativo) con media malla en $J$ hasta $40\,\tau_1$.}
\label{tab:colas}
\end{table}

Cuatro conclusiones. Primera: una distribución con $F_{\rm eq}(0)>0$ y
perturbación $s=1$ deja en $\delta\Phi$ una cola 150 veces mayor que la
gaussiana en $t_{\rm fin}$, en la frecuencia $\Omega_{\max}$; con
$s=(J/J_t)^{3/2}$ baja al nivel de la gaussiana. Segunda: esa perturbación lisa
produce poco potencial ($0.11$ de la gaussiana por unidad de masa), porque la
masa térmica está en órbitas casi circulares, de epiciclo chico, donde toda
perturbación lisa debe anularse; hay que compensarlo con $\varepsilon$ mayor
(sección~\ref{sec:perfil}). Tercera: la condición inicial solo controla la parte
\emph{libre}. La respuesta autoconsistente cerca de $J=0$ es
$\delta F_1\propto F_{\rm eq}'(0)\,\delta\Phi_1(J)$, con $\delta\Phi_1\propto\sqrt J$:
la autogravedad regenera una componente con la regularidad de $s=\sqrt{J/J_t}$ y
cola $t^{-2}$ en $\Omega_{\max}$, con amplitud proporcional al acoplamiento. En
una distribución que se anula como $J^2$, $F_{\rm eq}'(0)=0$ y esa componente no
aparece. Es una predicción, no una medida; su amplitud la da
\texttt{lineal.py} antes de cada corrida (sección~\ref{sec:protocolo}). Cuarta:
con cualquier familia la parte libre de $\delta\Phi$ vale todavía
$\sim10^{-4}$ de su valor inicial en $10\,\tau_1$, el orden del piso de ruido:
$\gamma$ se ajusta contra esa referencia, no suponiendo que la mezcla libre ya
terminó.

\subsection{El corrimiento de frecuencias por autogravedad}

La autogravedad hunde el potencial y acelera todas las órbitas: la banda entera
sube. Para entender cuánto, conviene mirar la órbita casi circular, que fija
$\Omega_{\max}$.

\paragraph{Aproximación 4: la frecuencia epicíclica.} Para una órbita cercana a la
circular de radio $r_c$ (definido por $\Phi'(r_c)=L_0^2/r_c^3$), la frecuencia
radial es la epicíclica,
\[
  \kappa^2 = \Phi_{\rm ef}''(r_c) = \Phi''(r_c) + \frac{3L_0^2}{r_c^4}
           = \Phi'' + \frac{3\Phi'}{r_c}.
\]
Usando la ecuación de Poisson en simetría esférica,
$\Phi''+2\Phi'/r = 4\pi G\rho$, se elimina $\Phi''$:
\[
  \boxed{\;\kappa^2 = 4\pi G\,\rho(r_c) + \frac{G\,M(<r_c)}{r_c^3}\;}
\]
La frecuencia radial depende de dos cosas: la masa encerrada, como en Kepler, y
\emph{la densidad local en el radio de la órbita}. Esta segunda es la clave.

\paragraph{Comprobación con el isócrono.} Con $L_0=2$, $r_c=5.742$. La densidad
del isócrono de masa y escala unidad es
$\rho = [3(1+a)a^2 - r^2(1+3a)]/[4\pi(1+a)^3a^3]$ con $a=\sqrt{1+r^2}$; en $r_c$
vale $1.09\times10^{-4}$, y la masa encerrada es $0.696$. Entonces
$4\pi\rho = 1.37\times10^{-3}$, $M/r_c^3 = 3.68\times10^{-3}$, y
$\kappa = 0.0711$, que coincide con $\Omega(0) = c^{-3} = 0.0711$. La fórmula
epicíclica es exacta en el límite circular, como debe.

\paragraph{La contribución propia.} La componente autogravitante de masa $a_0$
vive en una cáscara alrededor de $r_c$. El semiancho de la órbita con acción $J$
es la amplitud epicíclica, $A=\sqrt{2J/\kappa}$; tomando $J=\sigma_J$ (el máximo de
$J^2e^{-J^2/\sigma_J^2}$), $A=0.92$, $1.19$ y $1.68$ para $\sigma_J=0.03$, $0.05$
y $0.10$. El máximo de $J^2(J_t-J)^2$ está en $J_t/2=1.4\,\sigma_J$: una cáscara
parecida. La Maxwelliana de Wilson, en cambio, es máxima en $J=0$ y su acción
media es $0.20\,J_t=0.027$ para $J_t=0.138$: $A=0.87$ contra $1.19$, una cáscara
más angosta y más densa.

\paragraph{Aproximación 5: cáscara uniforme.} Si la masa se reparte uniforme en
$[r_c-A,\,r_c+A]$, su densidad es $\rho_{\rm self}\approx a_0/(4\pi r_c^2\,2A)$, y
la mitad de la masa queda dentro de $r_c$. El cambio de $\kappa^2$ es
\[
  \delta\kappa^2 \approx \frac{a_0}{2A\,r_c^2} + \frac{a_0/2}{r_c^3},
  \qquad
  \frac{\delta\kappa}{\kappa} = \frac{\delta\kappa^2}{2\kappa^2}.
\]

\begin{table}[htbp]
\centering
\begin{tabular}{lcccc}
\toprule
$\sigma_J$ & $4\pi\rho_{\rm self}/a_0$ & $M_{\rm self}/(r_c^3a_0)$ &
  $\delta\kappa/(\kappa a_0)$ estimado & $\delta\Omega/(\Omega a_0)$ medido\\
\midrule
0.10 & $9.0\times10^{-3}$ & $2.6\times10^{-3}$ & 1.15 & 1.55\\
0.05 & $1.28\times10^{-2}$ & $2.6\times10^{-3}$ & 1.52 & 1.87\\
0.03 & $1.65\times10^{-2}$ & $2.6\times10^{-3}$ & 1.89 & 2.19\\
\bottomrule
\end{tabular}
\caption{El corrimiento por unidad de masa. La estimación epicíclica con cáscara
uniforme reproduce el orden de magnitud, la tendencia (más corrimiento cuanto
más angosta la cáscara) y entre el \SI{74}{\%} y el \SI{86}{\%} del valor
medido; acierta más cuanto más angosta la cáscara, donde la hipótesis de
órbitas casi circulares es mejor. Subestima porque la distribución $J^2e^{-J^2/\sigma^2}$ tiene muchas
órbitas de amplitud chica y concentra la masa cerca de $r_c$ más que una cáscara
uniforme.}
\end{table}

Tres conclusiones de este cálculo, todas usadas en el diseño:

\begin{enumerate}
\item \textbf{El término de densidad local domina}: con $\sigma_J=0.05$ aporta el
\SI{83}{\%} de $\delta\kappa^2$. Por eso el corrimiento es más de dos veces lo que
predice la estimación ingenua por masa encerrada,
$\delta\Omega/\Omega\approx\tfrac12 a_0/M_{\rm iso}(r_c) = 0.72\,a_0$: la masa
propia está justo donde viven las órbitas.
\item \textbf{Angostar $\sigma_J$ ayuda dos veces}: la banda se estrecha
\emph{y} el corrimiento por unidad de masa crece (1.55, 1.87 y 2.19 para
$\sigma_J=0.10$, $0.05$ y $0.03$). Por eso $\eta$ crece más rápido que $1/\sigma_J$.
\item \textbf{El corrimiento es lineal en $a_0$} en todo el rango de la batería:
los valores calculados con \texttt{eta.py} dan coeficientes constantes al
\SI{5}{\%} entre $a_0=0.008$ y $0.125$.
\item \textbf{La forma de $F_{\rm eq}$ importa por la densidad local.} Con
$J^2(J_t-J)^2$ el coeficiente es $1.81$ y $2.08$ para $J_t=0.138$ y $0.083$,
cerca de la gaussiana. La Maxwelliana de Wilson, que concentra la masa en $r_c$,
da $2.15$, $2.58$ y $3.0$ para $J_t=0.276$, $0.138$ y $0.083$, y además
\emph{ensancha} la banda: el término $4\pi G\rho(r_c)$ sube $\Omega(0)=\kappa(r_c)$
mucho más que $\Omega(J_t)$, que corresponde a órbitas excéntricas que pasan
poco tiempo en la cáscara. Con $J_t=0.138$ el ancho relativo pasa de $0.162$ a
$0.194$ en $\eta=1$; en las familias huecas baja de $0.166$ a $0.136$. El mismo
$\eta$ corresponde así a geometrías distintas, y esa diferencia es lo que pone a
prueba el control E1.
\end{enumerate}

\subsection{El modo colectivo y el criterio $\eta$}

\paragraph{La respuesta lineal.} Se escribe la distribución como equilibrio más
perturbación, $F = F_{\rm eq}(J) + \delta F(Q,J,t)$, y el potencial igual,
$\Phi = \Phi_{\rm eq} + \delta\Phi$. A primer orden, la ecuación de Vlasov en
variables acción-ángulo es
\[
  \partial_t\,\delta F + \Omega(J)\,\partial_Q\,\delta F
  - \partial_Q\,\delta\Phi\;\frac{dF_{\rm eq}}{dJ} = 0 .
\]
Para un modo $\propto e^{i(kQ-\omega t)}$ se despeja
\[
  \delta F_k(J) = \frac{k\,F_{\rm eq}'(J)}{k\,\Omega(J)-\omega}\;\delta\Phi_k(J).
\]
El denominador se anula donde $k\Omega(J)=\omega$: son las partículas
\emph{resonantes}. La perturbación de la distribución produce una perturbación
de densidad, y esta, vía Poisson, un $\delta\Phi$; exigir que ese $\delta\Phi$ sea
el mismo que se supuso da una ecuación de autoconsistencia cuyas soluciones son
los modos. Esquemáticamente,
\[
  \varepsilon(\omega) = 1 - (\text{acoplamiento}\propto G\,a_0)\int dJ\;
  \frac{k\,F_{\rm eq}'(J)\,|\Psi(J)|^2}{k\,\Omega(J)-\omega} = 0 ,
\]
con $\Psi$ la forma radial del modo. En el caso real el acoplamiento es una
matriz en una base de funciones radiales (el método de Kalnajs), no un escalar;
la forma escalar basta para el argumento.

\paragraph{Dónde puede vivir el modo.} Si $\omega$ cae dentro de la banda, la
integral tiene un polo real: la raíz se desplaza al plano complejo con parte
imaginaria negativa y el modo se amortigua (amortiguamiento de Landau). Si
$\omega$ cae fuera, no hay polo, la raíz es real y el modo oscila sin
amortiguarse.

Para perturbaciones radiales de un sistema esférico hay una forma más precisa,
la de Antonov: la ecuación linealizada se escribe
$\partial_t^2\psi = -\mathcal A\psi$ con $\mathcal A$ autoadjunto y positivo, cuyo
espectro esencial es $\{k^2\Omega(J)^2\}$. El término gravitatorio de $\mathcal A$
es \emph{atractivo}, así que tiende a bajar los autovalores: un modo discreto,
si existe, aparece \emph{por debajo} de $\Omega_{\min}^2$, o en un hueco entre las
bandas de $k$ consecutivos. Esos huecos existen cuando
$\Omega_{\max}/\Omega_{\min}<2$, y aquí vale $1.37$ para $\sigma_J=0.10$.

\paragraph{Aproximación 6: $\delta\Omega$ como medida del acoplamiento.} El corrimiento
de la frecuencia media del equilibrio \emph{no es} la distancia que recorre el
modo; es un efecto del mismo acoplamiento sobre el estado de fondo. Lo usamos
como medida de la intensidad del acoplamiento porque ambos son del mismo orden:
los dos salen de términos $\sim 4\pi G\rho_{\rm self}$, como mostró la fórmula
epicíclica. La pregunta ``¿alcanza el acoplamiento para sacar el modo del
continuo?'' se convierte así en
\[
  \boxed{\;\eta = \frac{\delta\Omega}{\Delta\Omega}\;\gtrsim\;1\;?}
\]

\begin{nota}
\textbf{Lo que $\eta$ es y lo que no.} Es un criterio de diseño con fundamento
físico, no un teorema. El enunciado riguroso es dónde cae el autovalor de
$\mathcal A$ respecto del espectro esencial, y eso depende también de la
densidad de estados cerca del borde de la banda (sección~\ref{sec:borde}). Lo
que $\eta$ garantiza es que la batería explora el rango donde la transición
\emph{puede} ocurrir, en vez de quedarse entera de un solo lado. Quien decide es
la simulación.
\end{nota}

\subsection{Lo que ya dicen las corridas de Landau}
\label{sec:landau}

Las corridas regeneradas de la sección~11 de las notas tienen $a_0=10^{-2}$ y
$\sigma_J=0.10$, es decir $\eta\approx0.05$: muy adentro del continuo. Sus
resultados ilustran el marco y fijan la predicción para la batería.

\begin{itemize}
\item \textbf{La cola se amortigua, como corresponde a $\eta\ll1$.} El ajuste de la
cola da $\omega = 0.0570$ y $\gamma = 5.1\times10^{-3}$, estables en tres ventanas
temporales, con tres órdenes del ajuste y entre $N_{rc}=400$ y $800$.
\item \textbf{La frecuencia está dentro de la banda, cerca del borde inferior.} Con
$a_0=10^{-2}$ la banda es $[0.0522,\,0.0714]$: $\omega$ queda al \SI{25}{\%} de su
ancho desde abajo. Es lo que predice la forma de Antonov: el acoplamiento
atractivo empuja el polo hacia $\Omega_{\min}$, pero no alcanza para sacarlo.
\item \textbf{La autogravedad ya cambia la \emph{forma} del decaimiento.} Sin ella el
decaimiento es gaussiano (tabla anterior); con ella la cola es exponencial con
tasa $\gamma$. A tiempos largos la exponencial le gana a la gaussiana: el cociente
entre la señal con y sin autogravedad pasa de $1$ a $\sim30$ hacia $t=1600$.
\item \textbf{El régimen es lineal.} La corrida PIC coincide con la teoría lineal en
cuatro cifras ($0.41998$ contra $0.41994$ en $t=200$).
\item \textbf{El piso numérico limita la cola.} La corrida con $\varepsilon=0$ da el
nivel de ruido, $\approx10^{-4}$ en las mismas unidades. La cola se puede seguir
desde $t\approx1000$ hasta $t\approx1600$: unas tres $e$-foldings de $\gamma$. Es la
razón de ser de la corrida B2 con $N=10^5$.
\end{itemize}

\begin{nota}
\textbf{Predicción cuantitativa para la batería.} Al crecer $\eta$, la frecuencia de
la cola $\omega$ debe bajar hacia $\Omega_{\min}$ y la tasa $\gamma$ debe disminuir;
cuando $\omega$ quede por debajo de $\Omega_{\min}$, $\gamma$ debe ser cero dentro del
error. La curva $\gamma(\eta)$ y la posición de $\omega$ respecto de la banda son
el resultado central que la batería va a medir.
\end{nota}

\subsection{El eje no lineal: atrapamiento}
\label{sec:atrapamiento}

\paragraph{Por qué la no linealidad entra por $a_0\varepsilon$.} En un potencial fijo
la ecuación de Vlasov es \emph{exactamente lineal} en $F$: la solución de la
sección de mezcla de fases vale para cualquier $\varepsilon$. La no linealidad
aparece solo a través de $\delta\Phi$, que es proporcional a la masa y a la
amplitud de la perturbación.

\paragraph{Aproximación 7: el péndulo.} Cerca de una resonancia $k\Omega(J_r)=\omega$,
el hamiltoniano $H_0(J)+|\delta\Phi_k|\cos(kQ-\omega t)$ se reduce, en el sistema
que gira con la onda, a un péndulo: desarrollando $H_0$ a segundo orden en
$J-J_r$ se obtiene la frecuencia de rebote y el semiancho de la isla,
\[
  \omega_b = k\,\sqrt{|\Omega'(J_r)|\,|\delta\Phi_k|}, \qquad
  \Delta J_{\rm atr} = 2\sqrt{|\delta\Phi_k|/|\Omega'(J_r)|}.
\]
Supone una resonancia aislada y una amplitud constante durante un periodo de
rebote. El parámetro que compara atrapamiento con mezcla de fases es
\[
  \mu = \frac{\omega_b}{\Delta\Omega} .
\]
Si $\mu\ll1$ las partículas se desfasan antes de sentir el pozo y vale la teoría
lineal; si $\mu\gtrsim1$ quedan atrapadas antes, el amortiguamiento se detiene y
aparecen islas en el espacio fase. \textbf{Ninguna teoría actual cubre ese régimen}
para este sistema.

\paragraph{La amplitud de $\delta\Phi$: una corrección importante.} La estimación
ingenua es $|\delta\Phi|\sim\varepsilon\,|\Phi_{\rm self}|$, y es incorrecta por
un factor grande. Una perturbación $\cos Q$ agrupa las partículas en fase: hay
más cerca del pericentro que del apocentro. Pero eso no cambia la masa de la
cáscara, solo la \emph{desplaza} radialmente en una fracción $\sim\varepsilon A$
de su radio. El potencial de una cáscara desplazada difiere del original en
$\sim\varepsilon\,(A/r_c)\,|\Phi_{\rm self}|$, con $A/r_c\approx0.3$. La
corrida de Landau lo mide directamente: con $a_0=10^{-2}$ y $\varepsilon=0.1$,
\[
  |\delta\Phi|_{\max} = 2.3\times10^{-5}, \qquad
  |\delta\Phi|_{\rm rms}/|\Phi_{\rm self}| = 6.6\times10^{-3},
\]
unas quince veces menos que la estimación ingenua. Con $|\Omega'| = 0.0705$ y
$\Delta\Omega=0.0192$ resulta
\[
  \mu_{\rm Landau} = 0.066 \quad(0.045\ \text{con el valor rms}),
\]
consistente con la coincidencia en cuatro cifras con la teoría lineal. Esta es
la calibración del resto de la sección.

\paragraph{Aproximación 8: escalamiento de $\delta\Phi$.} De la imagen de la
cáscara desplazada, $|\delta\Phi|\propto a_0\,\varepsilon\,A\propto
a_0\,\varepsilon\,\sigma_J^{1/2}$ para la gaussiana; para el soporte compacto se
usa el ancho equivalente $\sigma_J=J_t/2.763$. La forma de $F_{\rm eq}$ y la
perturbación entran por la última columna de la tabla~\ref{tab:colas}, el
$\delta\Phi(0)$ por unidad de masa y de $\varepsilon$ relativo a la gaussiana,
$\rho_\Phi$: $1.09$ para $J^2(J_t-J)^2$ con $s=1$ y $0.11$ para Wilson con
$s=(J/J_t)^{3/2}$. Con la calibración anterior y el ancho de banda de cada caso,
\[
  \mu \approx 0.066\,
  \left(\frac{a_0\,\varepsilon}{10^{-3}}\right)^{1/2}
  \left(\frac{\sigma_J}{0.1}\right)^{1/4}
  \frac{0.0192}{\Delta\Omega}\,
  \left(\frac{|\Omega'|}{0.0705}\right)^{1/2}
  \rho_\Phi^{1/2},
\]
con $|\Omega'|=\Delta\Omega/J_t$ y $\rho_\Phi$ calculado con $J_t=0.138$ y
supuesto igual para los otros anchos.

\begin{table}[htbp]
\centering
\begin{tabular}{lllccc}
\toprule
Caso & $F_{\rm eq}$, ancho & $a_0$ & $\varepsilon$ & $\eta$ & $\mu$\\
\midrule
Landau (referencia) & gaussiana, $\sigma_J=0.10$ & 0.01 & 0.1 & 0.05 & 0.07\\
A1 & Wilson, $J_t=0.138$ & 0.0065 & 1.0 & 0.10 & 0.09\\
A4 & Wilson, $J_t=0.138$ & 0.075 & 1.0 & 1.00 & 0.25\\
A5 & Wilson, $J_t=0.138$ & 0.121 & 1.0 & 1.50 & 0.28\\
A9 & Wilson, $J_t=0.083$ & 0.044 & 1.0 & 1.00 & 0.26\\
\midrule
C4 & $J^2(J_t-J)^2$, $J_t=0.138$ & 0.075 & 1.0 & 1.00 & 0.97\\
C5 & $J^2(J_t-J)^2$, $J_t=0.083$ & 0.041 & 1.0 & 0.99 & 1.03\\
C6 & $J^2(J_t-J)^2$, $J_t=0.083$ & 0.100 & 1.0 & 3.1 & 1.7\\
\bottomrule
\end{tabular}
\caption{Valores estimados de $\mu$. Todo el bloque~A queda en régimen lineal
($\mu\le0.28$) aun con $\varepsilon=1$, porque la perturbación lisa de la
Maxwelliana produce poco potencial; por la misma razón el eje de atrapamiento se
recorre con la familia hueca (bloque~C). En ella, cruzar $\mu=1$ exige
$\varepsilon\approx1$, y con $\eta\approx1$ solo se llega justo al umbral. El valor de cada corrida se mide
a partir de su propio $\delta\Phi$.}
\label{tab:mu}
\end{table}

\begin{nota}
\textbf{Corrección respecto de la discusión previa.} En la conversación en que se
planteó este eje estimé $\mu$ con la amplitud ingenua de $\delta\Phi$, lo que daba
$\mu\approx1.75$ para $a_0=0.125$, $\varepsilon=1$ y $\sigma_J=0.10$. Con la
amplitud medida, ese caso da $\mu\approx0.75$. La conclusión cualitativa se
mantiene (el régimen atrapado es alcanzable con masas pequeñas), pero hace falta
$\varepsilon\approx1$ y un soporte angosto, y el bloque~C se rediseñó en
consecuencia.
\end{nota}

\paragraph{$\eta$ y $\mu$ no son independientes.} Con $\varepsilon$ fijo y
$\sigma_J$ el ancho (o su equivalente $J_t/2.763$), $\eta\propto a_0/\sigma_J$
(aproximadamente) y $\mu\propto a_0^{1/2}\sigma_J^{-3/4}$.
A $\eta$ fijo, $\mu\propto\eta^{1/2}\sigma_J^{-1/4}$ crece solo muy despacio al
angostar el soporte. La consecuencia es que la región ``modo sin amortiguar pero
todavía lineal'' ($\eta>1$, $\mu<1$) es accesible, mientras que la región
``atrapado pero sin modo'' ($\eta<1$, $\mu>1$) es prácticamente inaccesible con
esta familia: alcanzar $\mu=1.5$ con $\eta=0.5$ exigiría un soporte de ancho
$\sim10^{-3}$ en $J$, con tiempos de mezcla inviables. El mapa en $(\eta,\mu)$ tiene, en la práctica,
tres regiones y no cuatro.

\subsection{El eje del borde, y por qué no está en esta batería}
\label{sec:borde}

Hadžić, Rein, Schrecker y Straub (ARMA 2025, \texttt{arXiv:2301.07662}) prueban,
para equilibrios autogravitantes en el campo de una masa puntual central, una
dicotomía afilada: con $f\propto(E_0-E)^k$ en el borde del soporte, si $k>1$ las
perturbaciones lineales se amortiguan y si $1/2<k\le1$ no. La clave técnica de
la prueba es la ausencia de autovalores embebidos en el espectro esencial.

La sección de mezcla de fases da una intuición de por qué el borde importa: un
borde que se anula como $(\Omega-\Omega_{\min})^k$ produce en la respuesta libre una
cola algebraica $\propto t^{-(k+1)}$, más lenta cuanto menor es $k$. Con
autogravedad, una respuesta libre que decae despacio cerca del borde deja más
margen para que el acoplamiento sostenga una oscilación. Es una heurística, no
la prueba.

Hay que precisar qué bordes tiene la distribución de esta batería, la
Maxwelliana de Wilson. En $J=J_t$ (el extremo de baja frecuencia, $\Omega_{\min}$,
hacia donde el acoplamiento atractivo empuja al modo) tiene un borde de vacío
genuino, $F\propto(E_t-E)^2$, con exponente $k=g=2$: del lado ``amortigua'' de la
dicotomía y lejos del caso marginal $k=1$, que es King. En $J=0$ (el extremo de
alta frecuencia, $\Omega_{\max}$) es máxima, como toda distribución monótona y
como los equilibrios de ese trabajo. El mecanismo del borde no queda apagado
sino fijo en una clase conocida: cualquier oscilación que sobreviva debe
atribuirse a $\eta$. La ventaja frente a la gaussiana, que no tiene borde en
$\Omega_{\min}$, es que el experimento del borde pasa a ser esta misma batería
cambiando un solo número, $g$, a $\eta$ fijo (sección~\ref{sec:nocubre}).

\subsection{Resumen de aproximaciones}

\begin{table}[htbp]
\centering
\small
\begin{tabular}{p{0.8cm}p{4.3cm}p{3.4cm}p{5.1cm}}
\toprule
\# & Aproximación & Dónde se usa & Validación\\
\midrule
1 & Soporte $[0,J_t]$; umbral $10^{-2}F_{\max}$ solo en la gaussiana & bordes de banda, $\eta$ &
  exacto en la familia compacta; convención común en la gaussiana\\
2 & $\Delta\Omega/\Omega\approx3\Delta J/(J+c)$ & lectura de la banda &
  error $<\SI{1.3}{\%}$ contra el isócrono exacto\\
3 & $\tau_k=2\pi/(k\Delta\Omega)$ como escala & $t_{\rm fin}$ &
  decaimiento libre medido: gaussiano, $\Delta\Omega/\sigma_\Omega\approx6.2$--$7.0$\\
4 & Frecuencia epicíclica & origen del corrimiento &
  exacta en el límite circular: $\kappa=\Omega(0)=0.0711$\\
5 & Cáscara uniforme & estimación de $\delta\Omega$ &
  reproduce tendencia y \SI{74}{\%}--\SI{86}{\%} del valor; se usa el calculado\\
6 & $\delta\Omega$ mide el acoplamiento & definición de $\eta$ &
  heurística; la decide la simulación\\
7 & Péndulo (resonancia aislada) & definición de $\mu$ &
  estándar en dinámica no lineal; vale si $\mu$ no es $\gg1$\\
8 & $|\delta\Phi|\propto a_0\varepsilon\sigma_J^{1/2}\rho_\Phi$, $\sigma_J=J_t/2.763$ & tabla~\ref{tab:mu} &
  calibrada en una corrida; $\rho_\Phi$ calculado con $J_t=0.138$; factor $\sim2$\\
9 & Linealización de Vlasov & bloques A, B, D &
  PIC y teoría lineal coinciden en cuatro cifras con $\mu\approx0.07$\\
10 & Costo lineal en $N$ y en pasos & tiempos &
  calibrado con $N=10^4$; optimista para $N=10^5$ si escala peor en paralelo\\
11 & Colas libres en el isócrono desnudo ($a_0\to0$) & elección de $F_{\rm eq}$, referencia libre &
  cuadratura exacta; reproduce \texttt{landau\_\allowbreak libre}; exponentes confirmados\\
12 & Cola regenerada en $\Omega_{\max}$, $t^{-2}$, $\propto$ acoplamiento & lectura de la cola (Wilson) &
  predicción de la respuesta lineal; su amplitud la da \texttt{lineal.py}\\
\bottomrule
\end{tabular}
\caption{Aproximaciones del marco físico y del diseño, con cómo se comprobó
cada una.}
\label{tab:aprox}
\end{table}


\section{Decisiones de diseño y su justificación}

\subsection{La función de distribución}
\label{sec:familia}

Se usa la Maxwelliana rebajada de Wilson, escrita en la energía:
\[
  F_{\rm eq} = A\,E_\gamma\!\left(g,\,\frac{E_t-E}{T}\right) \ \ \text{si } E<E_t,
  \qquad
  E_\gamma(g,x) = e^x\,P(g,x) = \sum_{n\ge0}\frac{x^{g+n}}{\Gamma(g+n+1)},
\]
con $P$ la función gamma incompleta regularizada, $g=2$, $E_t=E(J_t)$ la energía
de la órbita de acción $J_t$ y $T=(E_t-E_c)/W_0$ con $E_c=E(0)$ la de la órbita
circular y $W_0=3$. $E_t$, $T$ y $A$ se recalculan en el potencial autoconsistente
de cada iteración, de modo que el soporte es siempre $[0,J_t]$ y la masa $a_0$. La
perturbación es $F=F_{\rm eq}[1+\varepsilon\,s(J)\,g(Q)]$ con $s=(J/J_t)^{3/2}$ y
$\varepsilon=1$ (sección~\ref{sec:perfil}).

\paragraph{Por qué no la gaussiana.} La gaussiana $J^2e^{-J^2/\sigma_J^2}$ se
propuso por su parecido con Maxwell--Boltzmann, pero no es una Maxwelliana. En la
distribución de rapideces de Maxwell, $4\pi v^2e^{-v^2/2\sigma^2}$, el $v^2$ es el
jacobiano de la cáscara de velocidades; la función de distribución misma es
$e^{-v^2/2\sigma^2}$, máxima en $v=0$. Aquí la medida ya es plana,
$dr\,dp_r=dQ\,dJ$, así que el $J^2$ de $F$ no es un jacobiano sino un vaciado
real del espacio fase alrededor de la órbita circular. En $r=r_c$, donde
$J\approx p_r^2/2\kappa$, da una distribución de $p_r$ proporcional a
$p_r^4e^{-p_r^4/4\kappa^2\sigma_J^2}$: nula en $p_r=0$ y con dos máximos, dos haces
radiales que se cruzan. Y tiene inversión de población, $\partial F/\partial E>0$
en $J<\sigma_J$.

\paragraph{El análogo correcto.} $F\propto e^{-E/T}$. Como
$E=p_r^2/2+\Phi_{\rm ef}(r)$, en cada radio da \emph{exactamente} una gaussiana en
$p_r$ con dispersión $\sqrt T$ y una densidad de Boltzmann
$e^{-\Phi_{\rm ef}/T}$, y como $E=E(J)$ es un equilibrio exacto. Cerca de la
órbita circular, $E\approx E_c+\kappa J$ y queda $F\propto e^{-\kappa J/\sigma_R^2}$,
la distribución cuasi-isoterma de Binney (2010, MNRAS 401, 2318), estándar para
discos tibios en variables de acción. El ancho tiene lectura física:
$J_0=\sigma_R^2/\kappa$, con $\sigma_R$ la dispersión de velocidades radiales; con
$W_0=3$ y $J_t=0.138$ resulta $\sigma_R/v_c=0.16$, una cáscara tibia.

\paragraph{El soporte compacto, por marea.} Un cúmulo en un campo externo pierde
las estrellas con $E>E_t$. La familia de isotermas rebajadas de Gieles y Zocchi
(2015, MNRAS 454, 576) describe ese corte: $g=1$ es King (1966), $e^x-1$; $g=2$ es
Wilson (1975), $e^x-1-x$. Cerca del borde $F\propto x^g/\Gamma(g+1)$: \textbf{$g$
es el exponente $k$ de Hadžić y colaboradores}; lejos del borde, $F\to e^x$. La
familia contiene como límites el politropo, $(E_t-E)^g$, con $W_0\to0$, y la
isoterma con $W_0\to\infty$.

\begin{table}[htbp]
\centering
\small
\begin{tabular}{L{4.4cm}L{2.9cm}L{2.9cm}L{2.9cm}}
\toprule
 & gaussiana & hueca & Wilson\\
 & $J^2e^{-J^2/\sigma_J^2}$ & $J^2(J_t-J)^2$ & $E_\gamma(2,(E_t-E)/T)$\\
\midrule
soporte compacto & no & sí & sí, por marea\\
banda & umbral del \SI{1}{\%} & exacta & exacta\\
distribución de $p_r$ en $r_c$ & bimodal & bimodal & gaussiana (rebajada)\\
en $J=0$ & $\propto J^2$ & $\propto J^2$ & máximo\\
monótona en $E$ (Antonov) & no & no & sí\\
perturbación lisa & $s=1$ & $s=1$ & $s=(J/J_t)^{3/2}$\\
cola libre de $\delta\Phi$ en $10\tau_1$ & $1.2\times10^{-4}$ & $4.2\times10^{-4}$ & $3.0\times10^{-4}$\\
$\delta\Phi(0)$ por unidad de masa, $\rho_\Phi$ & 1.00 & 1.09 & 0.11\\
$a_0$ para $\eta=1$ (ancho $0.138$) & 0.069 & 0.075 & 0.075\\
$\delta\Omega/(\Omega a_0)$ & 1.87 & 1.81 & 2.58\\
banda al agregar masa & se angosta & se angosta & se ensancha\\
\bottomrule
\end{tabular}
\caption{Las tres familias con el ancho de A4 ($\sigma_J=0.05$, $J_t=0.138$,
$L_0=2$). Colas y $\rho_\Phi$ de la tabla~\ref{tab:colas}; masas y coeficientes
de corrimiento de \texttt{eta.py}.}
\label{tab:familias}
\end{table}

Las razones de la elección, en orden de importancia:

\begin{enumerate}
\item \textbf{Es función solo de la energía, y por tanto de $J$}: un equilibrio
exacto del potencial total, que \texttt{equilibrio.py} construye iterando Poisson
hasta autoconsistencia. Si dependiera de $Q$ se enrollaría desde $t=0$, y ese
decaimiento sería indistinguible del que se quiere medir; sin equilibrio bien
definido no existe el marco en que se resta $\delta\Phi$.
\item \textbf{Es monótona}: $\partial F/\partial E<0$, así que por el teorema de
Antonov no tiene modos radiales crecientes. Una oscilación que no decae sería un
modo neutro genuino, sin la ambigüedad de una inestabilidad. Además, en la línea
de trabajos de Hadžić, Rein, Straub y colaboradores el operador linealizado se
define en un espacio con peso $1/|\varphi'(E)|$, lo que exige $\varphi'<0$ en todo
el soporte; si esto se confirma en la sección de hipótesis del artículo de
referencia, las dos familias huecas quedan fuera del teorema con que se quiere
comparar, y esta dentro.
\item \textbf{Su interior es térmico}: en cada radio la distribución de $p_r$ es
gaussiana, rebajada cerca de $E_t$; con $W_0=3$ el máximo vale \SI{80}{\%} del de
la isoterma sin truncar, $P(2,3)=0.80$. Sus parámetros tienen sentido físico:
temperatura, energía de marea y forma del borde.
\item \textbf{Su soporte compacto tiene causa física}, la marea, y deja la banda
exacta: el espectro esencial y $\eta$ son propiedades del equilibrio, no de un
umbral.
\item \textbf{Su borde es el de Hadžić y es ajustable}: $g=2$ está en la clase
``amortigua'', y el experimento del borde es esta misma batería con otro $g$,
incluido King ($g=1$), el caso marginal.
\end{enumerate}

\paragraph{Por qué $g=2$ y $W_0=3$.} Se exploró $g\in\{2,3\}$ y
$W_0\in\{1,\ldots,8\}$ con $J_t=0.138$ (tabla~\ref{tab:barrido}). Todas dan
$\eta=1$ con $a_0$ entre $0.066$ y $0.074$; lo que cambia es la forma. Subir $W_0$
hace más térmico el núcleo pero deja el borde sin peso y debilita la
perturbación lisa; con el mismo $W_0$, $g=2$ tiene el núcleo más térmico, la cola
más baja y más potencial que $g=3$. $W_0=3$ es el compromiso: núcleo al
\SI{80}{\%} de la isoterma y \SI{7}{\%} de la masa en la mitad exterior del
soporte, que es donde vive el borde. $W_0=6$, el límite casi térmico, entra como
control (E3).

\begin{table}[htbp]
\centering
\small
\begin{tabular}{llcccccc}
\toprule
$g$ & $W_0$ & $P(g,W_0)$ & $\langle J\rangle/J_t$ & masa en $J>J_t/2$ &
  $\sigma_R/v_c$ & cola libre $10\tau_1$ & $\rho_\Phi$\\
\midrule
2 & \textbf{3} & \textbf{0.80} & 0.20 & 0.068 & 0.157 & $3.0\times10^{-4}$ & 0.11\\
2 & 6 & 0.98 & 0.15 & 0.028 & 0.111 & $1.0\times10^{-3}$ & 0.06\\
3 & 3 & 0.58 & 0.17 & 0.037 & 0.157 & $6.8\times10^{-4}$ & 0.08\\
3 & 5 & 0.88 & 0.15 & 0.024 & 0.122 & $1.0\times10^{-3}$ & 0.06\\
\bottomrule
\end{tabular}
\caption{Candidatas de la familia de Wilson, con $J_t=0.138$, $L_0=2$ y
$s=(J/J_t)^{3/2}$. $P(g,W_0)$ es el cociente entre el máximo y el de la isoterma
sin truncar. El barrido completo en $W_0$ ($1$, $2$, $3$, $5$, $8$) muestra la misma
tendencia: $\langle J\rangle/J_t$ baja de $0.23$ a $0.12$ y la masa exterior de
\SI{10}{\%} a \SI{1}{\%}.}
\label{tab:barrido}
\end{table}

\paragraph{El costo: el borde en $\Omega_{\max}$.} Toda distribución monótona es
máxima en $J=0$, así que el borde superior de la banda, $\Omega_{\max}=\kappa(r_c)$,
tiene peso. Tres consecuencias, las tres medidas o calculadas:
\begin{enumerate}
\item[(a)] La perturbación $F_{\rm eq}(J)\cos Q$ es discontinua en la órbita
circular: allí $Q$ es el ángulo polar de unas coordenadas lisas $(x,y)$ con
$J\approx(x^2+y^2)/2$, y $F_{\rm eq}(0)\cos Q=F_{\rm eq}(0)\,x/|x|$. Su cola libre
en $\Omega_{\max}$ decae como $t^{-3/2}$, 150 veces por encima de la gaussiana en
$t_{\rm fin}$. Con $s=(J/J_t)^{3/2}$, que es lisa ($J^{3/2}\cos Q\propto Jx$), la
cola baja al nivel de la gaussiana (tabla~\ref{tab:colas}).
\item[(b)] Esa perturbación lisa produce poco potencial, $\rho_\Phi=0.11$, y exige
$\varepsilon=1$ para igualar el $\delta\Phi$ del diseño anterior. Es legítimo: la
no linealidad entra solo por $\delta\Phi$ (sección~\ref{sec:atrapamiento}), y
$\mu\le0.28$ en todo el bloque~A. La contracara es que la familia no alcanza el
régimen atrapado, que se estudia con la hueca (bloque~C).
\item[(c)] La autogravedad regenera una componente en $\Omega_{\max}$ con cola
$t^{-2}$ y amplitud proporcional al acoplamiento (sección de mezcla de fases),
que ninguna condición inicial elimina. No es un defecto de la distribución sino
del modelo de $L$ fijo: con $L$ distribuido, $\kappa(L)$ varía y difumina el borde.
Se ajusta como componente aparte, con frecuencia y exponente conocidos, y su
amplitud se calcula antes con \texttt{lineal.py}.
\end{enumerate}
Además la banda se ensancha con la masa (sección del corrimiento), y el mismo
$\eta$ describe otra geometría que en las familias huecas: el control E1 mide
si eso importa.

\paragraph{Por qué no el politropo puro.} $(E_t-E)^g$ es el límite $W_0\to0$ de esta
familia: conserva los tres costos anteriores y pierde el interior térmico.

\paragraph{La función de prueba.} $B(J)=J^2e^{-(J-j_1)^2/s_{j1}^2}$ sigue al
soporte: $j_1=s_{j1}=J_t/2.763$, la misma proporción que en 11\_landau ($0.10$
sobre un soporte que termina en $0.276$). Se fija en el \texttt{.par} de cada
corrida, de donde la leen \texttt{landau\_\allowbreak analisis.py} y, con
\texttt{--j1 --sj1}, \texttt{lineal.py}.

\subsection{El perfil angular de la perturbación}
\label{sec:perfil}

$g(Q)$ debe ser periódica y de media cero, que es lo único que se necesita para
que la masa no cambie y no haya que reconstruir el equilibrio. El generador
admite además un factor radial, $F=F_{\rm eq}\,[1+\varepsilon\,s(J)\,g(Q)]$:
$s=1$ (\texttt{--pert plana}), $s=\sqrt{J/J_{\max}}$ (\texttt{suave}) o
$s=(J/J_{\max})^{3/2}$ (\texttt{suave3}). Con $F_{\rm eq}(0)>0$ solo las dos últimas
son lisas en la órbita circular, y \texttt{suave3} es la que no deja cola en
$\Omega_{\max}$ (costo (a) de la sección~\ref{sec:familia}): es la de la familia
principal, con $\varepsilon=1$. La familia hueca y la gaussiana, que se anulan
como $J^2$, usan \texttt{plana} con $\varepsilon=0.1$, como 11\_landau. Como
$s\le1$, la cota de positividad es la misma.

\begin{itemize}
\item \texttt{cos}: $g=\cos Q$. Excita únicamente $k=1$. En régimen lineal cada
armónico evoluciona independiente, así que excitar uno solo da la medición más
limpia de $\gamma$ y $\omega$. Es el que se usa en toda la batería salvo donde
se indique.
\item \texttt{suma3}: $g=\cos Q+\tfrac12\cos 2Q+\tfrac13\cos 3Q$. Excita $k=1,2,3$
de entrada y sirve para la prueba de \emph{frecuencia común}, que es la firma de
un modo genuino. Su máximo es $\approx 1.5$, así que la positividad de $F$ exige
$\varepsilon\le 0.67$; el generador lo verifica y aborta si no se cumple.
\end{itemize}

\subsection{El momento angular $L_0$}

$L_0$ decide dónde vive el sistema dentro del isócrono, y con ello cuánta
autogravedad se obtiene por unidad de masa.

\begin{table}[htbp]
\centering
\begin{tabular}{rrrrr}
\toprule
$L_0$ & $c$ & radio circular $r_c$ & periodo radial & $M_{\rm iso}(r_c)$ \\
\midrule
0.25 & 1.133 & 0.804 & 9.1 & 0.078\\
0.5 & 1.281 & 1.30 & 13.2 & 0.192\\
1 & 1.618 & 2.42 & 26.6 & 0.413\\
\textbf{2} & \textbf{2.414} & \textbf{5.74} & \textbf{88.4} & \textbf{0.696}\\
3 & 3.303 & 10.9 & 226 & 0.828\\
\bottomrule
\end{tabular}
\caption{Con $L_0=2$ la distribución vive donde el fondo ya encierra el
\SI{70}{\%} de su masa, así que $a_0$ compite contra $0.7$. Con $L_0=1$ compite
contra $0.41$: la misma $\eta$ se alcanza con menos masa.}
\end{table}

El costo por corrida es casi independiente de $L_0$: el número de pasos para
cubrir el mismo número de tiempos de mezcla va como $T_r/r_c$, que vale $15.3$,
$11.0$ y $11.3$ para $L_0=2$, $1$ y $0.25$.

\textbf{Decisión:} $L_0=2$ como referencia, por continuidad con las corridas,
los equilibrios y los análisis ya validados; y dos corridas de control con
$L_0=1$. Si el resultado colapsa en $\eta$ con dos $L_0$, la afirmación deja de
depender de esa elección: $L_0$ cambia a la vez la masa de fondo encerrada, el
radio, el periodo y la forma de $\Omega(J)$.

Para comparar en igualdad de condiciones hay que escalar el ancho, porque lo que
importa es $\Delta\Omega/\Omega\approx 3\,\Delta J/c$: \textbf{el ancho del soporte
debe escalar como $c$}. El equivalente de $J_t=0.138$ con $L_0=2$ es
$J_t=0.094$ con $L_0=1$ ($\sigma_J=0.034$). Y la malla debe escalar con $r_c$: con $L_0=1$ el
sistema vive en $r\sim 1.5$--$4$, así que $\Delta r=0.05$ y $r_{\max}=8$ dan la
misma resolución relativa que $\Delta r=0.1$ y $r_{\max}=20$ con $L_0=2$.

\subsection{Por qué dos ejes y no solo la masa}

Con el ancho de las corridas de Landau ($J_t=0.276$, o $\sigma_J=0.10$) el
corrimiento es lineal en la masa, $\delta\Omega/\Omega\approx 2.1\,a_0$ con la
Maxwelliana, mientras que el ancho de la banda es $\approx 0.31$. Las masas
$10^{-3}$, $10^{-2}$ y $10^{-1}$ dan entonces $\eta\approx0.007$, $0.07$ y $0.68$:
\emph{las tres del mismo lado}. Un barrido así
mediría amortiguamiento en los tres casos y no distinguiría entre ``la
autogravedad no basta'' y ``no llegué al régimen''.

Con $J_t=0.083$ la banda sin masa baja a $0.10$ y la masa $a_0=0.044$ ya da
$\eta\approx 1.0$. Es decir: \textbf{el ancho del soporte es la perilla que
permite cruzar la transición sin salir de masas pequeñas}.

Se podría en cambio subir $a_0$ hasta $0.3$ o $0.5$ manteniendo
$J_t=0.276$, y también se cruzaría. No se hace porque: (i) mezcla dos
cambios, ya que subir la masa corre las frecuencias \emph{y} deforma el
equilibrio, de modo que no se puede atribuir la transición a una sola causa;
(ii) sale del régimen enunciado, porque con $a_0=0.5$ la componente
autogravitante es comparable a la masa de fondo encerrada; y (iii) aumenta el
riesgo de inestabilidad (sección~\ref{sec:riesgos}).

\subsection{Duración de las corridas y el techo de recurrencia}

Dos escalas mandan. La primera es la mezcla de fases: hace falta
$t_{\rm fin}\approx 10\,\tau_1$ para que la cola sea interpretable. La segunda
es la \emph{recurrencia} por aliasing de la malla de nodos en $J$, el artefacto
de Birdsall y Langdon visto desde el otro lado:
\[
  T_{\rm rec}(k) \;=\; \frac{2\pi}{k\,|d\Omega/dJ|\,\Delta J},
  \qquad \Delta J = \frac{J_{\max}}{N_{rc}} .
\]
Más allá de $T_{\rm rec}$ la meseta que se mida es un artefacto, y el límite es
\emph{distinto para cada $k$}: baja como $1/k$.

Con soporte compacto y $J_{\max}=J_t$ la cuenta se simplifica: $|d\Omega/dJ|\approx
\Delta\Omega/J_t$ y $\Delta J=J_t/N_{rc}$, así que
\[
  T_{\rm rec}(k) = \frac{N_{rc}}{k}\,\tau_1 ,
\]
independiente del ancho y de la masa. Con $N_{rc}=400$, $T_{\rm rec}(4)=100\,\tau_1
= 10\,t_{\rm fin}$ en toda la batería. Con la gaussiana (corrida E2),
$J_{\max}=6\sigma_J$ cubre 2.17 veces el soporte y el margen baja a
$4.5\,t_{\rm fin}$; por eso $J_{\max}$ debe seguir al soporte y no quedarse en su
valor histórico de $0.6$, que con $\sigma_J=0.03$ desperdiciaría la mayoría de
los nodos donde $F\approx0$. El generador usa por omisión $J_{\max}=J_t$ en la
familia compacta y $6\sigma_J$ en la gaussiana.

\subsection{Resolución}

Son cuatro fuentes de error independientes, y no se sustituyen entre sí:

\begin{itemize}
\item $N_{rc}$ (nodos en $J$): fija la ventana temporal antes de la recurrencia.
\item $N_{pc}$ (nodos en $Q$): fija cuántos modos $k$ son fiables; la
convergencia en $Q$ es espectral y $25$ resuelve bien hasta $k=4$.
\item $\Delta r$: error de la autogravedad (depósito y Poisson).
\item $\Delta t$: error del integrador temporal.
\end{itemize}

\section{La batería}

Configuración común salvo donde se indique: equilibrio autoconsistente con la
Maxwelliana de Wilson ($g=2$, $W_0=3$) y $J_{\max}=J_t$, $L_0=2$, malla
$r\in[0,20]$ con $\Delta r=0.1$, \texttt{courant}\,$=1$ (o sea $\Delta t=0.05$),
yoshida4, $N_{rc}=400$, $N_{pc}=25$ ($N=10^4$ partículas), $g(Q)=\cos Q$,
perturbación \texttt{suave3} con $\varepsilon=1$ y función de prueba
$j_1=s_{j1}=J_t/2.763$. Antes de cada corrida, \texttt{lineal.py} sobre el mismo
equilibrio da la referencia lineal (sección~\ref{sec:protocolo}).

\subsection{Bloque A: el eje $\eta$}

Los valores de $\eta$ de la tabla están \emph{calculados} con
\texttt{eta.py} sobre el equilibrio autoconsistente de cada caso, no estimados.

\begin{table}[htbp]
\centering
\begin{tabular}{llrrrrrr}
\toprule
\# & $J_t$ & $a_0$ & $\eta$ & banda $\Delta\Omega/\Omega$ & $\mu$ & $t_{\rm fin}$ & costo\\
\midrule
A1 & 0.138 & 0.0065 & 0.10 & 0.162 & 0.09 & 5500 & \SI{3.1}{min}\\
A2 & 0.138 & 0.020 & 0.30 & 0.169 & 0.15 & 5100 & \SI{2.9}{min}\\
A3 & 0.138 & 0.042 & 0.59 & 0.180 & 0.20 & 4600 & \SI{2.6}{min}\\
A4 & 0.138 & 0.075 & 1.00 & 0.194 & 0.25 & 3900 & \SI{2.2}{min}\\
A5 & 0.138 & 0.121 & 1.50 & 0.212 & 0.28 & 3300 & \SI{1.8}{min}\\
\midrule
A6 & 0.276 & 0.044 & 0.30 & 0.303 & 0.13 & 2900 & \SI{1.6}{min}\\
A7 & 0.276 & 0.148 & 1.00 & 0.319 & 0.22 & 2200 & \SI{1.2}{min}\\
A8 & 0.083 & 0.0109 & 0.30 & 0.107 & 0.15 & 8100 & \SI{4.5}{min}\\
A9 & 0.083 & 0.044 & 1.00 & 0.132 & 0.26 & 6000 & \SI{3.4}{min}\\
\bottomrule
\end{tabular}
\caption{A1--A5 barren $\eta$ con $J_t$ fijo. A6--A9 son la prueba de
degeneración: el mismo $\eta$ por caminos distintos. Si A2, A6 y A8 coinciden, y
A4, A7 y A9 coinciden, entonces manda $\eta$ y no la masa ni el ancho por
separado. $t_{\rm fin}=10\,\tau_1$ baja al subir $\eta$ porque la masa ensancha la
banda.}
\end{table}

\subsection{Bloque B: convergencia}

Todas sobre A4 ($\eta\approx 1$, el caso más delicado), cambiando un solo
parámetro por vez. Como $\Delta t=\texttt{courant}\,\Delta r/p_{\max}$, en B4 y B5
se compensa \texttt{courant} para que $\Delta t$ no cambie.

\begin{table}[htbp]
\centering
\begin{tabular}{llll}
\toprule
\# & Cambio & Qué prueba & Costo\\
\midrule
B1 & $N_{rc}=40$ ($N=10^3$) & resolución en $J$; $T_{\rm rec}(4)=t_{\rm fin}$ a propósito & \SI{0.2}{min}\\
B2 & $N_{rc}=4000$ ($N=10^5$) & lo mismo por arriba; fija el piso & \SI{22}{min}\\
B3 & $N_{pc}=50$ & modos $k$ fiables & \SI{4.4}{min}\\
B4 & $\Delta r=0.2$, \texttt{courant}\,$=0.5$ & error de la autogravedad & \SI{2.2}{min}\\
B5 & $\Delta r=0.05$, \texttt{courant}\,$=2$ & lo mismo por el otro lado & \SI{2.2}{min}\\
B6 & \texttt{courant}\,$=2$ & error del integrador & \SI{1.1}{min}\\
B7 & \texttt{courant}\,$=0.5$ & lo mismo por el otro lado & \SI{4.4}{min}\\
\bottomrule
\end{tabular}
\caption{Criterio de aceptación: $\gamma$ y $\omega$ medidos en la cola deben
moverse menos del \SI{5}{\%} al duplicar $N$, al partir $\Delta r$ a la mitad y
al partir $\Delta t$ a la mitad. Si alguno no cumple, esa configuración no se
reporta.}
\end{table}

\subsection{Bloque C: linealidad y atrapamiento}

Con la Maxwelliana y su perturbación lisa, $\mu\le0.28$ aun con $\varepsilon=1$:
la familia principal no llega al régimen atrapado (costo (b) de la
sección~\ref{sec:familia}). El bloque tiene por eso dos partes. C1 y C7 son sobre
A4: linealidad y frecuencia común en la familia principal. C2--C6 recorren el eje
$\mu$ con la familia hueca $J^2(J_t-J)^2$, que produce diez veces más potencial por
unidad de masa, a partir de E1 (bloque~E). El régimen atrapado es una pregunta
no lineal, sobre islas en el espacio fase, y no depende de la monotonía de
$F_{\rm eq}$ como la del modo discreto. Ninguna corrida requiere recalcular el
equilibrio, porque $g$ tiene media cero.

\begin{table}[htbp]
\centering
\begin{tabular}{lL{2.3cm}llccL{3.9cm}r}
\toprule
\# & base & $\varepsilon$ & $g(Q)$ & $\eta$ & $\mu$ & Qué prueba & costo\\
\midrule
C1 & A4 & 0.3 & \texttt{cos} & 1.00 & 0.13 & linealidad: $\gamma,\omega$ como en A4 & \SI{2.2}{min}\\
C2 & E1 & 0.3 & \texttt{cos} & 1.00 & 0.53 & primer paso hacia el atrapamiento & \SI{3.5}{min}\\
C3 & E1 & 0.6 & \texttt{cos} & 1.00 & 0.75 & atrapamiento incipiente & \SI{3.5}{min}\\
C4 & E1 & 1.0 & \texttt{cos} & 1.00 & 0.97 & al borde de $\mu=1$ & \SI{3.5}{min}\\
C5 & hueca, $J_t=0.083$, $a_0=0.041$ & 1.0 & \texttt{cos} & 0.99 & 1.03 & el punto $(\eta,\mu)\approx(1,1)$ & \SI{5.6}{min}\\
C6 & hueca, $J_t=0.083$, $a_0=0.10$ & 1.0 & \texttt{cos} & 3.1 & 1.7 & régimen atrapado claro & \SI{6.2}{min}\\
C7 & A4 & 0.6 & \texttt{suma3} & 1.00 & 0.19 & frecuencia común en $k=1,2,3$ & \SI{2.2}{min}\\
\bottomrule
\end{tabular}
\caption{$\mu$ estimado con la calibración de la corrida de Landau
(tabla~\ref{tab:mu}); el de cada corrida se mide de su propio $\delta\Phi$. C1
contra A4 es la prueba de linealidad de la familia principal; E1 contra C2, la
de la hueca. C6 es el único punto claramente del lado atrapado, y está a
$\eta\approx3$: como se explicó en la sección~\ref{sec:atrapamiento}, no se puede
tener $\mu>1$ sin $\eta\gtrsim1$. La suma de armónicos de C7 tiene máximo
$\approx1.5$, así que $\varepsilon=0.6$ está dentro del límite de positividad
($0.67$).}
\end{table}

\subsection{Bloque D: control con otro $L_0$}

\begin{table}[htbp]
\centering
\begin{tabular}{lllrrr}
\toprule
\# & $L_0$ & $J_t$ & $a_0$ & $\eta$ & $t_{\rm fin}$\\
\midrule
D1 & 1 & 0.094 & 0.018 & 0.29 & 1520\\
D2 & 1 & 0.094 & 0.069 & 1.00 & 1150\\
\bottomrule
\end{tabular}
\caption{Con malla adaptada al radio circular, $\Delta r=0.05$ y $r_{\max}=8$:
así $\Delta t$ se reduce a la mitad y cada corrida cuesta entre 1 y
\SI{2}{min}. Si D1 y D2 reproducen a A2 y A4, la afirmación queda demostrada
frente al cambio más agresivo posible.}
\end{table}

\subsection{Bloque E: controles de familia}

Tres corridas a $\eta\approx1$, cada una contra A4, para separar lo que depende
de $\eta$ de lo que depende de la forma de la distribución.

\begin{table}[htbp]
\centering
\small
\begin{tabular}{lL{2.9cm}rrlrlr}
\toprule
\# & $F_{\rm eq}$ & $a_0$ & $\eta$ & $s(J)$, $\varepsilon$ & $t_{\rm fin}$ & compara con & costo\\
\midrule
E1 & hueca, $J_t=0.138$ & 0.075 & 1.00 & 1, 0.1 & 6200 & A4 & \SI{3.5}{min}\\
E2 & gaussiana, $\sigma_J=0.05$ & 0.069 & 0.97 & 1, 0.1 & 6100 & E1 & \SI{3.4}{min}\\
E3 & Wilson, $W_0=6$ & 0.073 & 1.00 & $(J/J_t)^{3/2}$, 1 & 3500 & A4 & \SI{2.0}{min}\\
\bottomrule
\end{tabular}
\caption{E1 es $J^2(J_t-J)^2$: tiene el mismo borde de vacío ($k=2$) y el mismo
$\eta$ que A4, pero es hueca en $J=0$ y su banda se angosta con la masa en vez de ensancharse: si
coincide con A4, $\eta$ basta; si no, la geometría de la banda importa y $\eta$ debe
reformularse. Como no es monótona, un crecimiento en E1 se reporta como
inestabilidad. E2 es el puente con lo validado: la gaussiana de 11\_landau
($J_{\max}=0.3$), que debe coincidir con E1. E3 es el límite térmico de la
familia principal ($J_t=0.138$), con núcleo al \SI{98}{\%} de la isoterma: prueba que el resultado no depende de
$W_0$.}
\end{table}

\section{Protocolo de medida}
\label{sec:protocolo}

Antes de cada corrida PIC se resuelve el problema lineal sobre el mismo
equilibrio con \texttt{lineal.py} (minutos, sin partículas). Da $h_1$ y
$\delta\Phi$ lineales, incluida la componente regenerada en $\Omega_{\max}$ de la
familia principal, y es la referencia contra la que se lee la corrida: en
11\_landau la PIC coincidió con ella en cuatro cifras. De cada corrida se extraen
luego cuatro cosas; las dos primeras deciden.

\begin{enumerate}
\item \textbf{$\delta\Phi(r,t)$}, la perturbación del potencial respecto del
equilibrio, que calcula \texttt{landau\_\allowbreak analisis.py}. \emph{No depende de ningún
mapa acción-ángulo} y es lo que los teoremas dicen que decae o no. De su
envolvente tardía salen la tasa $\gamma$ y la frecuencia $\omega$. Su parte
libre no desaparece: decae algebraicamente ($\sim t^{-5/2}$,
tabla~\ref{tab:colas}) y en $10\,\tau_1$ vale todavía $\sim10^{-4}$ del valor
inicial. $\gamma$ y $\omega$ se ajustan donde la señal está claramente por
encima de esa referencia libre, que para cada corrida se calcula sin simular
(\texttt{colas\_\allowbreak libres.py}, \texttt{landau\_\allowbreak libre.py}). En la
familia principal el ajuste lleva además una componente en $\Omega_{\max}$ con
envolvente $t^{-2}$, la regenerada, con su amplitud tomada de la referencia
lineal.
\item \textbf{$h_k(t)$ para $k=0\ldots4$} en el marco del equilibrio. La firma de
un modo genuino es \emph{una frecuencia común a todos los $k$} con cocientes de
amplitud constantes. Si solo hay mezcla de fases, cada $k$ decae con su propia
tasa y no hay frecuencia común. Nótese que $h_0$ se conserva por construcción y
sirve de normalización, no de indicador dinámico.
\item \textbf{La banda $[\Omega_{\min},\Omega_{\max}]$ del equilibrio}, que se
calcula sin simular. Se compara $\omega$ con la banda: adentro significa
resonante y amortiguado; afuera, modo genuino. Se reporta el $\eta$ medido, no
el de diseño.
\item \textbf{Deriva de $J$ y conservación de la energía}, que acotan la ventana
temporal en la que el resultado es del sistema y no del método.
\end{enumerate}

\begin{table}[htbp]
\centering
\begin{tabular}{lll}
\toprule
Observación en la cola & Interpretación & Control que la confirma\\
\midrule
Decae, sin frecuencia común & mezcla de fases & sigue la referencia libre\\
Amplitud constante, $\omega$ común & modo genuino & $\omega$ fuera de la banda\\
Crece exponencialmente & inestabilidad & crece también con $\varepsilon$ menor\\
Meseta plana a nivel del piso & discretización & baja al subir $N$\\
\bottomrule
\end{tabular}
\caption{Cómo se lee cada resultado posible, y con qué control se distingue de
los demás.}
\end{table}

\section{Riesgos y controles}
\label{sec:riesgos}

\begin{enumerate}
\item \textbf{Inestabilidad en vez de oscilación.} La familia principal es
monótona y por Antonov no tiene modos radiales crecientes. La familia hueca (E1,
C2--C6) y la gaussiana (E2) sí crecen entre $J=0$ y su máximo, o sea
$\partial F/\partial E>0$ ahí: es justo la condición que puede sostener un modo
creciente. En ellas hay que ajustar la envolvente y confirmar amplitud
constante; si aparece crecimiento, se reporta como inestabilidad, no como ``no
amortigua''.
\item \textbf{Recurrencia.} Toda meseta medida más allá de $T_{\rm rec}(k)$ es
artefacto de la malla en $J$, y el límite baja con $k$. Al reportar mesetas por
modo hay que decir hasta qué $t$ vale cada una.
\item \textbf{Deriva numérica del equilibrio.} Si $J$ deriva apreciablemente
dentro de la ventana, lo que se mide ya no es el sistema lineal sino relajación
numérica. Ese es el límite real de la ventana útil, y puede ser más restrictivo
que la recurrencia.
\item \textbf{Piso de discretización.} Para sostener ``no amortigua'' hace falta
que el piso esté claramente por debajo de la amplitud observada durante varios
tiempos de mezcla. Es el papel de B2.
\item \textbf{Nodos casi circulares.} Con $J\to 0$ el mapa es degenerado, y en la
familia principal esos nodos son los de mayor peso. Las protecciones de
radicandos (punto E13 de la auditoría) cubren el caso, pero hay que verificar en
cada corrida que ninguna partícula se descarte.
\item \textbf{Cola de borde en la familia principal.} La componente de
$\delta\Phi$ en $\Omega_{\max}$ decae como $t^{-2}$; un ajuste de un solo término la
confundiría con un amortiguamiento lento. Está en el extremo de la banda opuesto
al del modo y se ajusta como componente aparte, con su frecuencia fija en el
borde y su amplitud de la referencia lineal.
\end{enumerate}

\section{Herramientas}

Los cambios ya están en el repositorio y son neutrales: con los valores por
omisión reproducen bit a bit lo anterior, lo que se verificó comparando el
estado inicial de Landau, la corrida lineal $400\times25$ y el análisis de una
corrida completa; tras agregar las familias compactas y la Maxwelliana se
repitieron las dos primeras, y también la de los estados iniciales de politropo
de la versión anterior. La única salida que cambió es la columna $T_{\rm rec}$
de \texttt{eta.py}: $|d\Omega/dJ|$ se tomaba derivando dos veces la tabla $E(J)$,
lo que daba ruido de orden uno entre nodos vecinos, y ahora es la pendiente
media sobre el soporte.

\begin{itemize}
\item \texttt{equilibrio.py}: acepta \texttt{--forma} \texttt{gauss} (con
\texttt{--sigma}), \texttt{politropo} ($J^m(J_t-J)^k$, con \texttt{--jt},
\texttt{--k}, \texttt{--m}) o \texttt{maxwell} ($E_\gamma(k,(E_t-E)/T)$, con
\texttt{--jt}, \texttt{--k}, \texttt{--w0}), \texttt{--jmax}, \texttt{--l0},
\texttt{--gq} y \texttt{--pert} \texttt{plana}/\texttt{suave}/\texttt{suave3};
$J_{\max}$ por omisión sigue al soporte ($J_t$ o $6\sigma_J$); verifica la
positividad de $F$ ante el $\varepsilon$ pedido; guarda en el archivo del
equilibrio $L_0$, el perfil angular, la forma, $(J_t,k,m,W_0)$ y la
perturbación. $E_\gamma$ se evalúa con su serie de términos positivos, exacta a
precisión de máquina.
\item \texttt{landau\_\allowbreak analisis.py}: lee $(j_1,s_{j1})$ de la función de prueba
del \texttt{params\_\allowbreak usados.par} de cada corrida (no son el $\sigma_J$ de la
distribución, sino parámetros del diagnóstico del código) y $L_0$ del
equilibrio.
\item \texttt{lineal.py}: toma del equilibrio $\sigma_J$, $J_{\max}$, $L_0$, la
forma y la perturbación (para la Maxwelliana, $dF/dJ=F'(E)\,\Omega(J)$ con la
tabla $E(J)$ del equilibrio); \texttt{--j1 --sj1} fijan la función de prueba.
\item \texttt{eta.py} (nuevo): dado $L_0$, la forma y $a_0$ construye el equilibrio
y reporta banda (exacta con soporte compacto), corrimiento, $\eta$, $\tau_1$,
$t_{\rm fin}$ sugerido, $T_{\rm rec}(k)$, número de pasos y costo estimado. Es el
guion con el que se eligieron las masas de este documento.
\item \texttt{colas\_\allowbreak libres.py} (nuevo): decaimiento libre exacto de
$h_1$ y $\delta\Phi$ por cuadratura, sin evolución temporal, validado contra
\texttt{landau\_\allowbreak libre}, y $\rho_\Phi$. Da la tabla~\ref{tab:colas}.
\end{itemize}

Ejemplo de uso:
\begin{verbatim}
python3 eta.py --forma maxwell --jt 0.138 --k 2 --w0 3 \
        --a0 0.0065 0.020 0.042 0.075 0.121
python3 equilibrio.py --forma maxwell --jt 0.138 --k 2 --w0 3 --a0 0.075 \
        --eps 1 --pert suave3 --nrc 400 --npc 25 --salida eta/A4.dat
python3 equilibrio.py --forma politropo --jt 0.138 --k 2 --m 2 --a0 0.075 \
        --eps 0.1 --nrc 400 --npc 25 --salida eta/E1.dat
python3 lineal.py eta/A4_equilibrio.npz eta/A4_lineal.npz --j1 0.05 --sj1 0.05
python3 colas_libres.py
\end{verbatim}

\section{Lo que este diseño no cubre}
\label{sec:nocubre}

\begin{itemize}
\item \textbf{La dicotomía del borde} (Hadžić y colaboradores): exige variar el
exponente del borde a $\eta$ fijo. Con la familia de esta batería es cambiar
\texttt{--k} (el $g$ de la Maxwelliana) a $\{0.75,1,1.5,3\}$ con el $J_t$, el $W_0$
y el $\eta$ de A4, sin tocar nada más; A4 ($g=2$) es el punto de partida y $g=1$
es exactamente King, el caso marginal.
\item \textbf{El momento angular distribuido.} El borde de $\Omega_{\max}$ y su
cola $t^{-2}$ son propios de $\delta(L-L_0)$; con una distribución en $L$,
$\kappa(L)$ varía y los difumina. Es la extensión natural del modelo, no de la
batería.
\item \textbf{El fondo de masa puntual} de ese trabajo. El código puede
reproducirlo exactamente (\texttt{BGtype = sphere} vale $-1/r$ fuera del radio
unidad), pero el mapa acción-ángulo numérico está escrito alrededor del
isócrono y habría que parametrizar el fondo.
\item \textbf{El régimen de masas comparables} ($a_0\gtrsim 0.3$), que es un
problema físico legítimo pero distinto, y con los riesgos de estabilidad ya
señalados.
\end{itemize}

\end{document}
